\documentclass[10pt,a4paper]{amsart}
\usepackage[margin=19mm]{geometry}
\usepackage{amsmath,amssymb,mathtools}
\usepackage[scr=rsfs,cal=euler]{mathalfa}
\usepackage{xfrac}
\usepackage[unicode,hidelinks,bookmarks=false]{hyperref}
\newcommand{\ie}{i.e.\ }
\newcommand{\cf}{cf.\ }
\newcommand{\resp}{respectively\ }
\newcommand{\Subsection}{\subsection}
\providecommand{\nobreakdash}{\nobreak\hbox{-}}
\setlength{\emergencystretch}{2em}
\multlinegap0pt
\usepackage{algorithm}\usepackage{algpseudocode}
\usepackage{algpseudocode}
\usepackage{amssymb}
\usepackage{bbm}
\usepackage{xcolor}
\usepackage{multirow}
\usepackage{tikz}
\usepackage{booktabs}
\newtheorem{theorem}{Theorem}
\newcommand{\rev}[1]{{#1}}
\title{Variable aggregation for nonlinear optimization problems}
\author{Sakshi Naik, Lorenz Biegler, Russell Bent, Robert Parker}
\date{Source: arXiv:2502.13869v3 (CC BY 4.0)}
\begin{document}
\maketitle

\noindent\emph{Source and licence.} Variable aggregation for nonlinear optimization problems. Version arXiv:2502.13869v3 (CC BY 4.0), distributed by arXiv under the Creative Commons Attribution 4.0 International licence, http://creativecommons.org/licenses/by/4.0/, as recorded in the arXiv licence metadata for this exact version and shown on the abstract page https://arxiv.org/abs/2502.13869v3. Full licence notice: \texttt{licenses/arxiv-2502.13869-CC-BY-4.0.txt}.\par\smallskip

\noindent\textit{Editorial scope.} Counted unit: the article's theorem thm:matching-bounds (source0.tex lines 744--754). The article's complete original proof of it is delivered with it (source0.tex lines 755--771, one \texttt{proof} environment of 17 lines). No mathematics is written, repaired or re-derived: the statement and the proof are the article's own lines, in the article's own notation, and no retained line is substituted or re-flowed.  Retained uncounted as the source-local context the proof needs: lines 701--724.  Nothing was omitted from any retained span, so the package carries no omission record.  No retained line contains a citation, so the wrapper carries no bibliography and no bibliography entry.  No retained line contains a figure environment, an image-inclusion command or drawing code, so no image is delivered and the package declares no asset dependency.  Editorial formatting only, in addition: an AMS \texttt{amsart} preamble that restates the article's own declarations and macros used by the retained lines, namely the environments \texttt{theorem} on separate counters, together with the standard packages those lines need.  The retained environments print in this document as Theorem 1 in order of appearance, whereas the article prints the same items with the numbering of its own full document.  The complete native source of the article as distributed in the arXiv e-print for this exact version is bundled unchanged as \texttt{sources/173547/source0.tex}.
The runtime of this algorithm is dominated by either the runtime of computing the
maximum linear matching $\rev{\mathcal{M}}_L$, or, if $G_L$ is small, the runtime of constructing
$G_L$ in the first place.
\begin{algorithm}
  \caption{: {\tt linear\_matching\_aggregation\_set}}
  \begin{algorithmic}[1]
    \State {\bf Inputs:} Sets of variables $X$ and constraints $C$
    \State $G = \text{\tt bipartite\_graph}(X, C)$
    \State $G_L = \text{\tt linear\_bipartite\_graph}(X, C)$
    \State $\mathcal{M}_L = \text{\tt maximum\_matching}(G_L)$
    \State $G_\mathcal{M} = \text{\tt induced\_subgraph}(G, \mathcal{M}_L)$
    \State $\mathcal{T} = \text{\tt block\_triangularize}(G_\mathcal{M}, \mathcal{M}_L)$
    \State $\mathcal{M}_T = \{\}$
    \For{$T\text{~\bf in } \mathcal{T}$}
      \If{$\text{\tt size}(T) = 1$}
        \State $\mathcal{M}_T \gets \mathcal{M}_T \cup T$
      \Else                                      
        \State $\mathcal{M}_T \gets \mathcal{M}_T \cup \text{\tt greedy}(T)$
      \EndIf
    \EndFor
    \State {\bf Return:} $\mathcal{M}_T$
  \end{algorithmic}
  \label{alg:matching}
\end{algorithm}

\begin{theorem}
  Let $X$ be a set of variables and $C$ a set of constraints, with bipartite graph $G$
  and linear bipartite graph $G_L$. Let $n_{\rm match}$ be the cardinality of $\mathcal{M}_L$,
  a maximum matching of the linear bipartite graph. Let $G_\mathcal{M}$ be the subgraph
  of $G$ induced by $\mathcal{M}_L$, with a block triangular partition $\mathcal{T}$
  of cardinality $n_{\rm block}$. Let $\mathcal{A}$ be the matching returned by Algorithm
  \ref{alg:matching} that is used for aggregation, with cardinality $n_{\rm agg}$.
  Then\rev{,}
  \[n_{\rm block} \leq n_{\rm agg} \leq n_{\rm match}\rev{.}\]
  \label{thm:matching-bounds}
\end{theorem}

\begin{proof}
  Each block $T\in \mathcal{T}$ contributes at least one edge to the matching $\mathcal{A}$.
  If $\texttt{size}(T)=1$, it contributes exactly one edge. If $\texttt{size}(T) > 1$, it
  contributes at least one edge as the first constraint encountered in
  {\tt greedy} contains at least one linear variable (the variable matched with
  this constraint in $\mathcal{M}_L$).
  Therefore, $n_{\rm block} \leq n_{\rm agg}$.

  Each block $T\in \mathcal{T}$ contributes at most ${\tt size}(T)$ edges to $\mathcal{A}$.
  This is by construction, as {\tt greedy} contributes at most one edge per
  constraint provided as input, and there are ${\tt size}(T)$ constraints provided as
  input. If ${\tt size}(T)=1$, then exactly ${\tt size}(T)$ edges are contributed to
  $\mathcal{A}$. The number of edges in $\mathcal{A}$ is therefore at most 
  $\sum_{T\in \mathcal{T}} {\tt size}(T)$, which equals $n_{\rm match}$ as $\mathcal{T}$
  partitions $\mathcal{M}_L$. That is,
  \[n_{\rm agg} \leq \sum_{T\in \mathcal{T}} {\tt size}(T) = n_{\rm match}\rev{.}\]
\end{proof}

\end{document}
